\documentclass[11pt]{article}

\usepackage[utf8]{inputenc}
\usepackage[T1]{fontenc}
\usepackage{lmodern}
\usepackage{geometry}
\usepackage{amsmath,amssymb,amsfonts,amsthm,mathtools}
\usepackage{bm}
\usepackage{enumitem}
\usepackage{microtype}
\usepackage{hyperref}
\usepackage{titlesec}
\usepackage{booktabs}
\usepackage{array}
\usepackage{setspace}
\usepackage{mathrsfs}
\usepackage{physics}

\geometry{margin=1in}
\hypersetup{colorlinks=true, linkcolor=black, urlcolor=blue, citecolor=black}
\setlist[enumerate]{topsep=0.4em,itemsep=0.35em}
\setlist[itemize]{topsep=0.3em,itemsep=0.25em}
\titleformat{\section}{\large\bfseries}{\thesection.}{0.5em}{}
\titleformat{\subsection}{\normalsize\bfseries}{\thesubsection.}{0.5em}{}
\setstretch{1.06}

\newcommand{\Aether}{\AE{}ther}
\newcommand{\Aflow}{\AE{}ther-flow}
\newcommand{\TheoryName}{\Aflow{} Interpretation of Relativity}
\newcommand{\TheoryProgram}{\Aether{} / \Aflow{} framework}
\newcommand{\Sub}{\mathcal{A}}
\newcommand{\BridgeMetric}{\mathfrak{g}}
\newcommand{\AY}{\mathcal A_Y}
\newcommand{\AZ}{\mathcal A_Z}
\newcommand{\FieldSpace}{\mathscr F}
\newcommand{\GaugeGroup}{\mathscr G_{\mathrm{of}}}

\newtheorem{definition}{Definition}
\newtheorem{proposition}{Proposition}
\newtheorem{remark}{Remark}
\newtheorem{corollary}{Corollary}
\newtheorem{theorem}{Theorem}

\title{\textbf{The \TheoryName{}}\\[0.4em]
\large Explicit Symmetry/Quotient Candidate for Ordered-Flow Label Redundancy}
\author{}
\date{}

\begin{document}

\maketitle

\begin{abstract}
The first honest post-exhaustion non-algebraic momentum-completion candidate has now been tested and found coefficient-level negative: it preserves a genuine ordered-flow equation, but its flat-background observer-equation correction tensor vanishes only by collapsing back to the already exhausted algebraic route. The next fixed-route burden is therefore not another momentum-completion note. It is the explicit symmetry/quotient route already identified in the tracking layer.

This note writes that route in the smallest honest form presently available while keeping the bridge metric and the single-metric matter route fixed. Start from the doubly tuned exact-square completion with
\[
\AY
\equiv
\frac{\sqrt{\lambda_4}}{M_*}\Delta_\perp S-M_YY,
\qquad
\AZ
\equiv
m_S\bigl(U^A\nabla_A S-\sigma_0\bigr)-M_ZZ,
\]
and impose the local ordered-flow relabeling symmetry
\[
S\mapsto S+\eta,
\qquad
Y\mapsto Y+\frac{\sqrt{\lambda_4}}{M_YM_*}\Delta_\perp\eta,
\qquad
Z\mapsto Z+\frac{m_S}{M_Z}U^A\nabla_A\eta,
\]
with \(G_{AB}\), \(U^A\), \(Q^I\), and the bridge metric held fixed. The full tuned action and the observer map then depend only on the gauge-invariant combinations \(\AY\), \(\AZ\), and \(\BridgeMetric_{AB}\). In that explicit quotient interpretation, the previously recorded infinite-dimensional substrate fiber is no longer a physical underdetermination; it is the gauge orbit of the ordered-flow label.

The note then performs the same immediate stopping test required everywhere else in the active sequence. On the admissible flat background, the scalar quadratic action depends only on the gauge-invariant linearized combinations
\[
\AY^{(1)}
=
\frac{\sqrt{\lambda_4}}{M_*}\Xi-M_Yy,
\qquad
\AZ^{(1)}
=
m_S\Psi_h-M_Zz,
\]
so the quadratic reduced scalar bridge channel vanishes identically after algebraic elimination:
\[
\mathcal L_{\mathrm{bridge,quot}}^{(2)}\equiv 0.
\]
The explicit flat-background observer-equation correction tensor is therefore
\[
\mathcal C_{\mu\nu}^{\mathrm{quot,br}}\equiv 0.
\]

This is a coefficient-level positive structural candidate, not a completed derivational verdict. The note stops exactly at the required point: the ordered-flow label has been made physically redundant by an explicit local quotient principle while the bridge metric and single-metric matter route remain fixed, and the flat-background observer-tensor test is passed on that quotient branch. No reduced-system, local-coercive, or theorem work is reopened here. If a later architecture screen were to show that this quotient interpretation cannot be implemented honestly beyond the present flat-background coefficient level, then the next larger move would be to reconsider the bridge metric or the matter route.
\end{abstract}

\section{Introduction}

The double exact-square replacement candidate and its architecture-obstruction follow-up left the active derivational program in a very specific state. The coefficient-level observer remainder could be removed on the doubly tuned branch, but only at the price of eliminating the ordered-flow equation and leaving an infinite-dimensional substrate fiber over the same observer solution. The broader fixed-route momentum-compensator note then showed that no finite algebraic/gapped enlargement of that least-drift mechanism does better. The first honest non-algebraic post-exhaustion candidate changed that picture only partly: it preserved a genuine ordered-flow equation, but it still failed the flat-background observer-tensor test unless it collapsed back to the exhausted algebraic route.

The remaining fixed-route direction is therefore the one already identified in the Markdown tracking layer: if one wants to keep the bridge metric and the single-metric matter route fixed, the next note must be an explicit symmetry/quotient note rather than another momentum-completion note. That is the task here.

\paragraph{Framework and claim boundary.}
\TheoryProgram{} denotes the broader \Aether{} / \Aflow{} framework built on the claims that the \Aether{} is the underlying four-dimensional substrate of reality, the \Aflow{} is its intrinsic ordered motion, observed three-dimensional space is the local experiential slice of that deeper substrate, S-time is the experienced order of change arising from matter, light, and the \Aflow{}, and observed expansion is the three-dimensional appearance of deeper four-dimensional ordered motion. Within that framework, \TheoryName{} denotes the exact-closure theory in which the effective gravitational dynamics are adopted to be exactly Einsteinian with universal matter coupling. In the active sequence, \emph{adoption} means use of that established relativistic dynamics without claiming substrate derivation, while \emph{derivation} is reserved for a first-principles recovery from explicit substrate variables.

The present note stays inside that boundary. It does not claim a completed microphysical derivation. It does not reopen reduced-system or theorem work. It performs the next structural candidate-selection task only: make the ordered-flow label physically redundant through an explicit local symmetry/quotient principle if that can be done honestly while keeping the bridge metric and single-metric matter route fixed, then stop again at flat-background quadratic reduction plus the explicit observer-equation correction tensor.

\section{What Counts as an Honest Symmetry/Quotient Route}

\begin{definition}[Honest ordered-flow quotient principle]
\label{def:honest}
An \emph{honest ordered-flow quotient principle} for a fixed-route candidate branch consists of:
\begin{enumerate}
    \item an explicit local field transformation group \(\GaugeGroup\) acting on the ordered-flow label variables while leaving the bridge metric and matter route unchanged;
    \item exact invariance of the candidate action and observer map under that transformation;
    \item a physical field space defined as the quotient \(\FieldSpace/\GaugeGroup\);
    \item the property that the previously observed substrate fiber over a fixed observer solution is reinterpreted as a gauge orbit rather than as a genuine physical underdetermination.
\end{enumerate}
\end{definition}

\begin{remark}
Definition \ref{def:honest} is stronger than simply noticing that many substrate representatives project to the same observer metric. A true quotient principle has to identify those representatives as the same physical state by an explicit local symmetry, not merely leave them unexplained.
\end{remark}

\section{Explicit Ordered-Flow Relabeling Symmetry}

\begin{definition}[Tuned exact-square quotient candidate]
\label{def:candidate}
The \emph{tuned exact-square quotient candidate} is the action
\begin{equation}
S_{4YZ}^{\mathrm{quot}}
=
\frac{c^3}{16\pi G_N}
\int_{\Sub} d^4X\,\sqrt{G}\,\bigl(R[\BridgeMetric]-2\Lambda_{\Sub}\bigr)
+\int_{\Sub} d^4X\,\sqrt{G}\,\mathcal L_{\mathrm{aux},4YZ}
+S_{\mathrm{matter}}[\Xi^*\BridgeMetric,\psi],
\label{eq:S4YZ}
\end{equation}
with bridge metric
\begin{equation}
\BridgeMetric_{AB}=G_{AB}-2U_AU_B,
\qquad
P_{AB}=G_{AB}-U_AU_B,
\label{eq:bridge}
\end{equation}
and auxiliary Lagrangian
\begin{align}
\mathcal L_{\mathrm{aux},4YZ}
=\;&
\lambda_U\bigl(G_{AB}U^AU^B-1\bigr)
-\frac12
\left(
\frac{\sqrt{\lambda_4}}{M_*}\Delta_\perp S-M_YY
\right)^2
\nonumber\\
&-\frac12
\left(
m_S\bigl(U^A\nabla_A S-\sigma_0\bigr)-M_ZZ
\right)^2
-\frac12\delta_{IJ}P^{AB}(D_AQ^I)(D_BQ^J)
\nonumber\\
&-\frac12\widehat M_{IJ}Q^IQ^J
-\frac{\kappa_\theta}{2}\bigl(\nabla_AU^A\bigr)^2
-\frac{\kappa_\Omega}{4}\Omega_{AB}\Omega^{AB},
\label{eq:L4YZ}
\end{align}
where
\begin{equation}
\Delta_\perp S\equiv \nabla_A(P^{AB}\nabla_BS),
\qquad
\Omega_{AB}\equiv P_A{}^CP_B{}^D\nabla_{[C}U_{D]},
\label{eq:defops}
\end{equation}
and
\begin{equation}
m_S^2>0,
\qquad
\lambda_4>0,
\qquad
M_Y^2>0,
\qquad
M_Z^2>0,
\qquad
\widehat M_{IJ}>0,
\qquad
\kappa_\theta>0,
\qquad
\kappa_\Omega>0.
\label{eq:positivity}
\end{equation}
\end{definition}

\begin{definition}[Ordered-flow relabeling group]
\label{def:gauge}
Let \(\GaugeGroup\) denote the local transformation group generated by smooth relabeling parameters \(\eta\) that preserve the chosen asymptotic ordered-flow profile. Its action on the fields of Definition \ref{def:candidate} is
\begin{equation}
\delta_\eta S=\eta,
\qquad
\delta_\eta Y=\frac{\sqrt{\lambda_4}}{M_YM_*}\Delta_\perp\eta,
\qquad
\delta_\eta Z=\frac{m_S}{M_Z}U^A\nabla_A\eta,
\label{eq:gauge}
\end{equation}
with
\begin{equation}
\delta_\eta G_{AB}=0,
\qquad
\delta_\eta U^A=0,
\qquad
\delta_\eta Q^I=0,
\qquad
\delta_\eta \lambda_U=0.
\label{eq:gaugefixed}
\end{equation}
\end{definition}

\begin{proposition}[Gauge-invariant exact-square combinations]
\label{prop:invariants}
Under Definition \ref{def:gauge}, the combinations
\begin{equation}
\AY
\equiv
\frac{\sqrt{\lambda_4}}{M_*}\Delta_\perp S-M_YY,
\qquad
\AZ
\equiv
m_S\bigl(U^A\nabla_A S-\sigma_0\bigr)-M_ZZ
\label{eq:AYAZ}
\end{equation}
are exactly invariant. So are the bridge metric \(\BridgeMetric_{AB}\), the observer metric \(g_{\mu\nu}=\Xi^*\BridgeMetric_{AB}\), and the matter route \(S_{\mathrm{matter}}[\Xi^*\BridgeMetric,\psi]\).
\end{proposition}

\begin{proof}
Apply \eqref{eq:gauge}. Since \(G_{AB}\) and \(U^A\) are fixed under the symmetry, one has
\[
\delta_\eta\!\left(\Delta_\perp S\right)=\Delta_\perp\eta,
\qquad
\delta_\eta\!\left(U^A\nabla_A S\right)=U^A\nabla_A\eta.
\]
Therefore
\[
\delta_\eta \AY
=
\frac{\sqrt{\lambda_4}}{M_*}\Delta_\perp\eta
-M_Y
\left(
\frac{\sqrt{\lambda_4}}{M_YM_*}\Delta_\perp\eta
\right)
=0,
\]
and similarly
\[
\delta_\eta \AZ
=
m_S U^A\nabla_A\eta
-M_Z
\left(
\frac{m_S}{M_Z}U^A\nabla_A\eta
\right)
=0.
\]
Because \(\BridgeMetric_{AB}=G_{AB}-2U_AU_B\) depends only on \(G_{AB}\) and \(U^A\), it is invariant under \eqref{eq:gaugefixed}; hence \(g_{\mu\nu}\) and the matter route are invariant as well.
\end{proof}

\begin{theorem}[Exact symmetry and quotient interpretation]
\label{thm:symmetry}
The candidate action \(S_{4YZ}^{\mathrm{quot}}\) of Definition \ref{def:candidate} is exactly invariant under the ordered-flow relabeling group \(\GaugeGroup\). Consequently the physical field space may be defined as the quotient \(\FieldSpace/\GaugeGroup\), and on that quotient the ordered-flow label is physically redundant rather than merely underdetermined.
\end{theorem}

\begin{proof}
By Proposition \ref{prop:invariants}, the bridge metric, the matter action, and the two exact-square combinations \(\AY\) and \(\AZ\) are invariant under \(\GaugeGroup\). Every term in \eqref{eq:S4YZ}--\eqref{eq:L4YZ} depends on the fields only through \(\BridgeMetric_{AB}\), \(\AY\), \(\AZ\), \(Q^I\), \(U^A\), and \(\lambda_U\), all of which are invariant. Hence the action is exactly invariant.

Since the observer map depends only on \(\BridgeMetric_{AB}\), it is constant on \(\GaugeGroup\)-orbits. The same is true of the matter route. It is therefore consistent, at the present candidate-selection scope, to define the physical field space as the quotient \(\FieldSpace/\GaugeGroup\). On that quotient the ordered-flow label \(S\) is not an independent physical observable; it is a gauge representative.
\end{proof}

\begin{corollary}[The old infinite-dimensional fiber becomes gauge]
\label{cor:fiber}
On the exact-square locus
\begin{equation}
\AY=0,
\qquad
\AZ=0,
\label{eq:locus}
\end{equation}
the infinite-dimensional substrate fiber recorded in the old architecture-obstruction note coincides with the \(\GaugeGroup\)-orbit through any representative.
\end{corollary}

\begin{proof}
If \((S_1,Y_1,Z_1)\) and \((S_2,Y_2,Z_2)\) obey \eqref{eq:locus} for the same \(G_{AB}\), \(U^A\), and \(Q^I\), define \(\eta=S_2-S_1\). Then
\[
M_Y(Y_2-Y_1)
=
\frac{\sqrt{\lambda_4}}{M_*}\Delta_\perp(S_2-S_1)
=
\frac{\sqrt{\lambda_4}}{M_*}\Delta_\perp\eta,
\]
and similarly
\[
M_Z(Z_2-Z_1)
=
m_S U^A\nabla_A(S_2-S_1)
=
m_S U^A\nabla_A\eta.
\]
Thus \((S_2,Y_2,Z_2)\) is obtained from \((S_1,Y_1,Z_1)\) by the transformation \eqref{eq:gauge}. The previously recorded fiber is therefore the gauge orbit once the quotient principle is imposed explicitly.
\end{proof}

\begin{remark}
Corollary \ref{cor:fiber} is the structural point of the present manuscript. Nothing has been changed in the bridge metric or the matter route. The new input is that the old multiplicity of ordered-flow labels is now identified as local gauge redundancy.
\end{remark}

\section{Flat-Background Quadratic Reduction}

Use the same admissible homogeneous flat background as in the previous tuned and replacement notes:
\begin{equation}
\bar G_{AB}=\delta_{AB},
\qquad
\bar U^A=\delta^A{}_0,
\qquad
\bar S=\sigma_0X^0,
\qquad
\bar Q^I=0,
\qquad
\bar Y=0,
\qquad
\bar Z=0.
\label{eq:background}
\end{equation}
Write perturbations as
\begin{equation}
G_{AB}=\delta_{AB}+H_{AB},
\qquad
U^A=\bar U^A+V^A,
\qquad
S=\sigma_0X^0+\sigma,
\qquad
Y=y,
\qquad
Z=z,
\label{eq:perturb}
\end{equation}
with scalar metric variables
\begin{equation}
h_{00}=-2\phi=-H_{00},
\qquad
h_{0i}=S_i+\partial_iB=-H_{0i}-2V_i,
\qquad
V_i=V_i^T+\partial_i\chi.
\label{eq:metricpert}
\end{equation}

\begin{proposition}[Linearized ordered-flow relabeling]
\label{prop:linGauge}
Around the flat background \eqref{eq:background}, the linearized symmetry generated by \(\eta\) is
\begin{equation}
\delta_\eta \sigma=\eta,
\qquad
\delta_\eta y=\frac{\sqrt{\lambda_4}}{M_YM_*}\nabla^2\eta,
\qquad
\delta_\eta z=\frac{m_S}{M_Z}\partial_0\eta,
\label{eq:lingauge}
\end{equation}
with \(\delta_\eta h_{\mu\nu}=0\), \(\delta_\eta q^I=0\), and \(\delta_\eta\chi=0\).
\end{proposition}

\begin{proof}
Linearize \eqref{eq:gauge} on the background \eqref{eq:background}. Since \(\bar U^A=\delta^A{}_0\) and \(\bar P^{ij}=\delta^{ij}\), one has
\[
\delta_\eta\!\left(\Delta_\perp S\right)^{(1)}=\nabla^2\eta,
\qquad
\delta_\eta\!\left(U^A\nabla_A S\right)^{(1)}=\partial_0\eta,
\]
which gives \eqref{eq:lingauge}.
\end{proof}

\begin{proposition}[Gauge-invariant linearized exact-square blocks]
\label{prop:lininv}
Define
\begin{equation}
\Psi_h\equiv \partial_0\sigma-\sigma_0\phi,
\qquad
\Xi\equiv \nabla^2\sigma+\sigma_0\nabla^2(B+\chi),
\label{eq:psixi}
\end{equation}
and
\begin{equation}
\AY^{(1)}
=
\frac{\sqrt{\lambda_4}}{M_*}\Xi-M_Yy,
\qquad
\AZ^{(1)}
=
m_S\Psi_h-M_Zz.
\label{eq:linearAYAZ}
\end{equation}
Then \(\AY^{(1)}\) and \(\AZ^{(1)}\) are invariant under the linearized symmetry \eqref{eq:lingauge}.
\end{proposition}

\begin{proof}
Using Proposition \ref{prop:linGauge},
\[
\delta_\eta \Xi=\nabla^2\eta,
\qquad
\delta_\eta \Psi_h=\partial_0\eta.
\]
Substituting into \eqref{eq:linearAYAZ} gives \(\delta_\eta\AY^{(1)}=0\) and \(\delta_\eta\AZ^{(1)}=0\).
\end{proof}

\begin{proposition}[Quadratic scalar action on the quotient candidate]
\label{prop:quad}
At fixed bridge perturbation \(h_{\mu\nu}\), the complete flat-background quadratic nonmetric scalar action of Definition \ref{def:candidate} is
\begin{align}
\mathcal L_{\mathrm{aux},4YZ}^{(2)}
=\;&
-\frac12 q^I\!\left((-\nabla^2)\delta_{IJ}+\widehat M_{IJ}\right)\!q^J
-\frac{\kappa_\theta}{2}
\left(
\nabla^2\chi+\frac12\partial_0H
\right)^2
\nonumber\\
&-\frac12\bigl(\AY^{(1)}\bigr)^2
-\frac12\bigl(\AZ^{(1)}\bigr)^2.
\label{eq:quad}
\end{align}
\end{proposition}

\begin{proof}
Expand \eqref{eq:L4YZ} to quadratic order around \eqref{eq:background}. The \(Q\)- and \(\chi\)-blocks are the same as in the previous tuned notes. The two tuned scalar squares reduce exactly to \(-\frac12(\AY^{(1)})^2-\frac12(\AZ^{(1)})^2\). There is no other quadratic scalar term.
\end{proof}

\begin{remark}
Proposition \ref{prop:quad} makes the quotient structure explicit at coefficient level: the scalar quadratic action depends on the ordered-flow variables only through the gauge-invariant combinations \(\AY^{(1)}\) and \(\AZ^{(1)}\).
\end{remark}

\begin{proposition}[Quadratic reduced scalar bridge channel]
\label{prop:reduction}
After algebraic elimination of \(y\) and \(z\), and after the same decaying-branch longitudinal solve
\begin{equation}
\nabla^2\chi=-\frac12\partial_0H,
\label{eq:chisolve}
\end{equation}
the quadratic reduced scalar bridge channel of the quotient candidate vanishes identically:
\begin{equation}
\mathcal L_{\mathrm{bridge,quot}}^{(2)}\equiv 0.
\label{eq:zero}
\end{equation}
\end{proposition}

\begin{proof}
Varying \eqref{eq:quad} with respect to \(y\) and \(z\) gives the algebraic equations
\begin{equation}
M_Yy=\frac{\sqrt{\lambda_4}}{M_*}\Xi,
\qquad
M_Zz=m_S\Psi_h,
\label{eq:yzsolve}
\end{equation}
equivalently
\[
\AY^{(1)}=0,
\qquad
\AZ^{(1)}=0.
\]
Substituting \eqref{eq:yzsolve} back into \eqref{eq:quad} removes both exact-square scalar blocks identically. The remaining \(\chi\)-equation is exactly \eqref{eq:chisolve}, as in the previous tuned notes, and it removes the old longitudinal auxiliary piece in the same way. No scalar bridge term remains. Hence \eqref{eq:zero} holds.
\end{proof}

\begin{corollary}[Gauge fixing of the ordered-flow perturbation]
\label{cor:gaugefix}
On the decaying perturbative branch one may choose the gauge \(\sigma=0\) by taking \(\eta=-\sigma\). In that gauge the vanishing \eqref{eq:zero} is manifest, but the result is gauge independent because \(\mathcal L_{\mathrm{bridge,quot}}^{(2)}\) was derived from the invariant quantities \(\AY^{(1)}\) and \(\AZ^{(1)}\).
\end{corollary}

\begin{proof}
The gauge choice \(\sigma=0\) is allowed by Proposition \ref{prop:linGauge} for decaying \(\sigma\). Gauge independence follows from Proposition \ref{prop:lininv}.
\end{proof}

\section{Explicit Observer-Equation Correction Tensor}

\begin{definition}[Quotient-branch correction tensor]
\label{def:Cquot}
At flat-background quadratic scope, define the \emph{quotient-branch observer-equation correction tensor} by
\begin{equation}
\mathcal C_{\mu\nu}^{\mathrm{quot,br}}
=
8\pi G_N\,T_{\mu\nu}^{\mathrm{quot}},
\label{eq:Cquot}
\end{equation}
where
\begin{equation}
T_{\mu\nu}^{\mathrm{quot}}
=
\rho_{\mathrm{quot}}\,n_\mu n_\nu
+2n_{(\mu}J_{\nu)}^{\mathrm{quot}}
+S_{\mu\nu}^{\mathrm{quot}},
\label{eq:Tquot}
\end{equation}
with \(J_\mu^{\mathrm{quot}}n^\mu=0\) and \(S_{\mu\nu}^{\mathrm{quot}}n^\nu=0\).
\end{definition}

\begin{theorem}[Flat-background observer-tensor verdict]
\label{thm:verdict}
For the quotient candidate of Definition \ref{def:candidate},
\begin{equation}
\rho_{\mathrm{quot}}=0,
\qquad
J_i^{\mathrm{quot}}=0,
\qquad
S_{ij}^{\mathrm{quot}}=0,
\label{eq:components}
\end{equation}
and therefore
\begin{equation}
T_{\mu\nu}^{\mathrm{quot}}\equiv 0,
\qquad
\mathcal C_{\mu\nu}^{\mathrm{quot,br}}\equiv 0.
\label{eq:verdict}
\end{equation}
\end{theorem}

\begin{proof}
By Proposition \ref{prop:reduction}, the reduced scalar bridge channel vanishes identically at the flat-background quadratic level. Hence there is no surviving scalar energy density, momentum density, or spatial stress contribution from that channel. Equations \eqref{eq:components} follow, and then \eqref{eq:verdict} is immediate from Definition \ref{def:Cquot}.
\end{proof}

\begin{remark}
Theorem \ref{thm:verdict} is exactly the stopping test required by the tracking layer. At the present coefficient level, the explicit symmetry/quotient route does remove the flat-background observer-equation correction tensor while keeping the bridge metric and single-metric matter route fixed.
\end{remark}

\section{Current Status and Stopping Point}

\begin{corollary}[Current coefficient-level status]
\label{cor:status}
At the present disciplined scope, the explicit symmetry/quotient route is an honest structural candidate and is coefficient-level positive on the flat admissible background.
\end{corollary}

\begin{proof}
Theorem \ref{thm:symmetry} establishes that the ordered-flow label is physically redundant through an explicit local quotient principle rather than through mere underdetermination. Theorem \ref{thm:verdict} establishes that the same quotient candidate has vanishing flat-background observer-equation correction tensor while preserving the bridge metric and matter route. This is precisely the coefficient-level candidate-selection burden assigned to the present note.
\end{proof}

\begin{corollary}[What this note does not reopen]
\label{cor:stop}
The present manuscript does not by itself reopen reduced-system, local-coercive, or theorem work on the quotient candidate.
\end{corollary}

\begin{proof}
The active tracking instructions require any new structural candidate to stop at flat-background quadratic reduction plus its explicit observer-equation correction tensor before broader nonlinear work is revisited. The present note has done exactly that and no more.
\end{proof}

\begin{remark}
If a later architecture screen were to show that the quotient principle of Definition \ref{def:gauge} cannot be implemented honestly beyond the current flat-background coefficient level, then the next larger move would be to reconsider the bridge metric or the single-metric matter route. That contingency is recorded here, but it is not the current verdict.
\end{remark}

\section{Conclusion}

The present manuscript records the next live structural step after the failed non-algebraic momentum-completion candidate.
\begin{enumerate}
    \item It writes an explicit local ordered-flow relabeling symmetry on the tuned double exact-square branch while keeping the bridge metric and the single-metric matter route fixed.
    \item It shows that the previously recorded infinite-dimensional substrate fiber is thereby reinterpreted as a gauge orbit, so the ordered-flow label is physically redundant rather than merely underdetermined.
    \item It carries that quotient candidate through the same required flat-background quadratic reduction and finds that the scalar bridge channel vanishes identically after algebraic elimination.
    \item It computes the corresponding explicit observer-equation correction tensor and finds \(\mathcal C_{\mu\nu}^{\mathrm{quot,br}}\equiv0\) at that coefficient level.
\end{enumerate}

That is the honest stopping point requested by the tracking layer. The explicit symmetry/quotient route now exists as a concrete manuscript-level candidate. It passes the flat-background observer-tensor test while leaving the bridge metric and single-metric matter route fixed. No reduced-system or theorem work is reopened here.

\end{document}
